\documentclass[10pt, aspectratio=169]{beamer}

% --- PACKAGES ---
\usepackage[utf8]{utf8}
\usepackage[vietnamese]{babel}
\usepackage{amsmath, amsfonts, amssymb}
\usepackage{booktabs}
\usepackage{graphicx}
\usepackage{tikz}

% --- THEME & STYLING (CLEAN ACADEMIC BEAMER) ---
\usetheme{default}
\usecolortheme{seahorse}

\setbeamertemplate{navigation symbols}{}
\setbeamertemplate{footline}[frame number]

\definecolor{academicnavy}{RGB}{15, 23, 42}
\definecolor{academicblue}{RGB}{67, 56, 202}

\setbeamercolor{frametitle}{fg=academicnavy, bg=white}
\setbeamercolor{title}{fg=academicnavy}
\setbeamercolor{structure}{fg=academicblue}

% --- TITLE METADATA ---
\title{\textbf{HỆ THỐNG AI DỰ ĐOÁN KHÁCH HÀNG QUAY LẠI MUA SẮM}}
\subtitle{Customer Repurchase Prediction — Academic ML Report}
\author{Kỹ sư AI \& Full-Stack Data Science DMC}
\institute{DMC AI Laboratory --- E-Commerce Analytics}
\date{Tháng 9, 2026}

\begin{document}

% === SLIDE 1: TITLE ===
\begin{frame}
    \titlepage
\end{frame}

% === SLIDE 2: PROBLEM FORMULATION ===
\begin{frame}{\textbf{1. Tổng Quan Bài Toán \& Phát Biểu Toán Học}}
    \begin{block}{\textbf{Bối Cảnh Nghiệp Vụ E-Commerce}}
        Trong E-Commerce, chi phí thu hút một khách hàng mới ($\text{CAC}$) đắt gấp $5 \div 7$ lần chi phí giữ chân khách hàng hiện hữu ($\text{Retention}$). Việc nhận diện sớm rủi ro rời bỏ ($\text{Churn Risk}$) giúp doanh nghiệp phân bổ chính xác ngân sách khuyến mãi.
    \end{block}

    \vspace{0.2cm}
    \begin{block}{\textbf{Phát Biểu Toán Học Bài Toán Học Máy}}
        Cho tập dữ liệu huấn luyện $\mathcal{D} = \{(\mathbf{x}_i, y_i)\}_{i=1}^N$ với $N = 5,630$ mẫu quan sát:
        \begin{itemize}
            \item $\mathbf{x}_i \in \mathbb{R}^8$: Vector $8$ thuộc tính RFM+ đại diện cho lịch sử hành vi.
            \item $y_i \in \{0, 1\}$: Nhãn nhị phân ($y_i = 1$ nếu phát sinh đơn mới trong 30 ngày; $y_i = 0$ nếu rời bỏ Churn).
        \end{itemize}
        \textbf{Mục tiêu:} Huấn luyện hàm phân loại $f(\mathbf{x}): \mathbb{R}^8 \to [0, 1]$ ước lượng xác suất hậu phương $P(y = 1 \mid \mathbf{x})$.
    \end{block}
\end{frame}

% === SLIDE 3: DATASET & FEATURES ===
\begin{frame}{\textbf{2. Bộ Dữ Liệu Kaggle \& Bảng Đặc Trưng RFM+}}
    \begin{columns}
        \column{0.48\textwidth}
        \textbf{Thống Kê Tập Dữ Liệu Kaggle:}
        \begin{itemize}
            \item Tổng số mẫu quan sát: $N = 5,630$ khách hàng.
            \item Phân bố nhãn mục tiêu (\texttt{will\_return}):
            \begin{itemize}
                \item Nhãn $1$ (Quay lại): $4,684$ mẫu ($83.2\%$)
                \item Nhãn $0$ (Churn): $946$ mẫu ($16.8\%$)
            \end{itemize}
        \end{itemize}

        \column{0.5\textwidth}
        \textbf{Danh Sách 8 Thuộc Tính Vector $\mathbf{x}$:}
        \begin{table}
            \small
            \begin{tabular}{ll}
                \toprule
                \textbf{Ký hiệu} & \textbf{Ý nghĩa Nghiệp vụ} \\
                \midrule
                \texttt{age} & Độ tuổi khách hàng (18--70) \\
                \texttt{gender} & Giới tính (Nam/Nữ) \\
                \texttt{total\_purchases} & Tần suất mua hàng ($F$) \\
                \texttt{avg\_order\_value} & Chi tiêu TB/đơn ($M$) \\
                \texttt{days\_since\_last\_purchase} & Số ngày chưa mua ($R$) \\
                \texttt{membership\_level} & Hạng thành viên \\
                \texttt{used\_voucher} & Lịch sử sử dụng voucher \\
                \texttt{satisfaction\_score} & Điểm hài lòng ($1 \div 5$) \\
                \bottomrule
            \end{tabular}
        \end{table}
    \end{columns}
\end{frame}

% === SLIDE 4: PIPELINE & ANTI-LEAKAGE ===
\begin{frame}{\textbf{3. Sơ Đồ Pipeline Processing \& Nguyên Tắc Chống Rò Rỉ Dữ Liệu}}
    \textbf{Sơ Đồ Quy Trình Pipeline:}
    \begin{center}
        \small
        \texttt{[Raw Kaggle Data (5,630)]} $\to$ \texttt{[Preprocessing \& Scaler]} $\to$ \texttt{[80/20 Stratified Split]} $\to$ \texttt{[GridSearchCV 5-Fold]}
    \end{center}

    \vspace{0.3cm}
    \begin{alertblock}{\textbf{Nguyên Tắc Phòng Chống Rò Rỉ Dữ Liệu (Anti-Data Leakage Protocol)}}
        Tất cả các đại lượng thống kê trung bình $\mu$ và độ lệch chuẩn $\sigma$ của StandardScaler cùng quy tắc One-Hot Encoder được thực hiện \textbf{chỉ trên tập Train (fit)}. Tập Test độc lập chỉ áp dụng phép biến đổi \textbf{transform}, đảm bảo tuyệt đối không có sự rò rỉ tri thức từ tập kiểm thử.
    \end{alertblock}
\end{frame}

% === SLIDE 5: MODEL BENCHMARK ===
\begin{frame}{\textbf{4. Kết Quả Đánh Giá \& So Sánh 3 Thuật Toán Học Máy}}
    \begin{table}
        \centering
        \caption{Kết quả đánh giá thực nghiệm trên tập Test độc lập ($N_{\text{test}} = 1,126$)}
        \begin{tabular}{lccccc}
            \toprule
            \textbf{Thuật toán} & \textbf{Accuracy} & \textbf{Precision} & \textbf{Recall} & \textbf{F1-Score} & \textbf{ROC-AUC} \\
            \midrule
            \textbf{Random Forest (Chính)} & \textbf{88.10\%} & \textbf{0.9051} & \textbf{0.9573} & \textbf{0.9304} & \textbf{0.9241} \\
            Decision Tree & 87.30\% & 0.9324 & 0.9135 & 0.9228 & 0.7936 \\
            Logistic Regression & 67.50\% & 0.9398 & 0.6506 & 0.7689 & 0.7845 \\
            \bottomrule
        \end{tabular}
    \end{table}

    \vspace{0.2cm}
    \begin{itemize}
        \item \textbf{Recall 95.73\%:} Giúp hạn chế tối đa việc bỏ sót khách hàng tiềm năng quay lại mua sắm.
        \item \textbf{ROC-AUC 0.9241:} Phản ánh năng lực phân tách cực tốt giữa 2 nhóm nhãn (Quay lại vs Churn).
    \end{itemize}
\end{frame}

% === SLIDE 6: CONFUSION MATRIX & FEATURE IMPORTANCE ===
\begin{frame}{\textbf{5. Sơ Đồ Ma Trận Nhầm Lẫn \& Trọng Số 8 Đặc Trưng}}
    \begin{columns}
        \column{0.48\textwidth}
        \textbf{Ma trận Nhầm lẫn (Confusion Matrix):}
        \begin{table}
            \small
            \begin{tabular}{ccc}
                \toprule
                & \textbf{Dự đoán Churn (0)} & \textbf{Dự đoán Return (1)} \\
                \midrule
                \textbf{Thực tế Churn (0)} & $TN = 96$ & $FP = 94$ \\
                \textbf{Thực tế Return (1)} & $FN = 40$ & $TP = 896$ \\
                \bottomrule
            \end{tabular}
        \end{table}

        \column{0.5\textwidth}
        \textbf{Top Trọng Số Đặc Trưng (Gini Importance):}
        \begin{enumerate}
            \item \texttt{age} (Độ tuổi): $32.68\%$
            \item \texttt{avg\_order\_value} (AOV): $28.57\%$
            \item \texttt{days\_since\_last\_purchase}: $12.62\%$
            \item \texttt{satisfaction\_score}: $8.53\%$
            \item \texttt{total\_purchases}: $7.93\%$
        \end{enumerate}
    \end{columns}
\end{frame}

% === SLIDE 7: WEB APP ARCHITECTURE ===
\begin{frame}{\textbf{6. Sơ Đồ Kiến Trúc Hệ Thống Nền Tảng Web App}}
    \textbf{Sơ Đồ Kiến Trúc End-to-End:}
    \begin{center}
        \small
        \texttt{[Frontend Dashboard]} $\longleftrightarrow$ \texttt{[Flask RESTful API]} $\longleftrightarrow$ \texttt{[ML Inference Engine (.joblib)]}
    \end{center}

    \vspace{0.3cm}
    \begin{itemize}
        \item \textbf{Giao diện Minimal White Luxury:} Loại bỏ hoàn toàn icon thừa, làm nổi bật trực quan con số xác suất phần trăm $\%$ tái mua.
        \item \textbf{3 Phân hệ Admin:} Real-time Predictor, Benchmark Evaluation Dashboard, Batch CSV Prediction Module.
    \end{itemize}
\end{frame}

% === SLIDE 8: CONCLUSION & ROADMAP ===
\begin{frame}{\textbf{7. Kết Luận Báo Cáo \& Định Hướng Phát Triển}}
    \begin{block}{\textbf{Kết Luận Giá Trị Kinh Doanh}}
        \begin{itemize}
            \item Tăng $15 \div 20\%$ Tỷ lệ Giữ chân Khách hàng ($\text{Retention Rate}$) và tiết kiệm $30\%$ Chi phí Voucher Marketing.
        \end{itemize}
    \end{block}

    \vspace{0.2cm}
    \begin{block}{\textbf{Định Hướng Nâng Cấp Kỹ Thuật Tương Lai}}
        \begin{itemize}
            \item \textbf{Survival Analysis (Phân Tích Sinh Tồn):} Áp dụng Cox Proportional Hazards Model để dự đoán chính xác số ngày trước khi khách phát sinh đơn mới.
            \item \textbf{MLOps \& Automated Retraining:} Tự động hóa pipeline huấn luyện lại mô hình định kỳ và giám sát hiện tượng Data Drift.
        \end{itemize}
    \end{block}
\end{frame}

\end{document}
